\section{The morphism Hall framework in type $A$}
\label{sec:prior}

This section fixes conventions and recalls the exact-category results used by
the new quotient construction.  We include the statements needed later so
that the Hall--skein argument can be read independently of the earlier
iterations.

\subsection{The exact morphism category}

Let $A$ be a basic finite-dimensional algebra over a field $\kk$.  Write
$\proj A$ for the category of finitely generated projective right $A$-modules
and
\[
 \MM=\Mor(\proj A).
\]
An object of $\MM$ is a map
\[
 f:P^{-1}\longrightarrow P^0,
\]
and a morphism from $f$ to $g:Q^{-1}\to Q^0$ is a commutative square.  The
exact structure is degreewise split: a sequence is a conflation when its
sequences in degrees $-1$ and $0$ split in $\proj A$.  The category is
Krull--Schmidt, finitary over a finite field, and has homological dimension at
most one \cite{Bautista,Li}.

For a module $N$, fix a minimal projective presentation
\[
 P_1^N\xrightarrow{d_N}P_0^N\longrightarrow N\longrightarrow0
\]
and put
\[
 C_N=(P_1^N\xrightarrow{d_N}P_0^N),
 \qquad
 Z_P=(P\longrightarrow0),
 \qquad
 U_P=(0\longrightarrow P),
 \qquad
 K_P=(P\xrightarrow{1_P}P).
\]
The $K_P$ are contractible and projective-injective in the two-term homotopy
window.

The following normal form was proved in the preceding paper by splitting off
invertible matrix blocks and then using minimality.

\begin{theorem}[Morphism homology normal form]
\label{thm:normal-form}
For every $f:P^{-1}\to P^0$ there are projectives $P_f,Q_f$, unique up to
isomorphism, such that
\begin{equation}
 f\simeq C_N\oplus Z_{P_f}\oplus K_{Q_f},
 \qquad N=H^0(f)=\coker f.
 \label{eq:normal-form}
\end{equation}
Moreover
\begin{equation}
 H^{-1}(f)=\ker f\simeq\Omega_A^2N\oplus P_f.
 \label{eq:homology-residue}
\end{equation}
Consequently the full homology pair recovers the reduced object
$f_{\rm red}=C_N\oplus Z_{P_f}$ through Krull--Schmidt cancellation
\[
 P_f=H^{-1}(f)\ominus\Omega_A^2H^0(f).
\]
Its split index is
\begin{equation}
 I(f)=[P^0]-[P^{-1}]
     =g_A(N)-[P_f]
 \quad\text{in }K_0^{\rm sp}(\proj A).
 \label{eq:split-index}
\end{equation}
\end{theorem}

The notation $\ominus$ in the theorem means subtraction of actual
Krull--Schmidt multiplicities.  It is not a formal difference in a Grothendieck
group.

The extension space in $\MM$ is computed by the two-term Hom complex.  For
$X=(P^-\xrightarrow fP^0)$ and $Y=(Q^-\xrightarrow gQ^0)$ there is a natural
exact sequence
\begin{equation}
\begin{aligned}
0\longrightarrow\Hom_{\MM}(X,Y)
&\longrightarrow
 \Hom_A(P^-,Q^-)\oplus\Hom_A(P^0,Q^0)\\
&\xrightarrow{(a,b)\mapsto ga-bf}
 \Hom_A(P^-,Q^0)
 \longrightarrow\Ext^1_{\MM}(X,Y)\longrightarrow0.
\end{aligned}
\label{eq:Hom-Ext-complex}
\end{equation}
For reduced objects the formula may be expressed in module-theoretic form.

\begin{proposition}[Reduced extension formula]
\label{prop:reduced-Ext}
Let $X=C_M\oplus Z_P$ and $Y=C_N\oplus Z_Q$.  Then
\begin{align}
 \Ext^1_{\MM}(X,Y)
 &\simeq D\Hom_A(N,\tau M)\oplus\Hom_A(P,N),
 \label{eq:Ext-XY}\\
 \Ext^1_{\MM}(Y,X)
 &\simeq D\Hom_A(M,\tau N)\oplus\Hom_A(Q,M).
 \label{eq:Ext-YX}
\end{align}
In particular, full morphism homology determines rigidity and pairwise
compatibility.
\end{proposition}

The correspondence of Li identifies support $\tau$-rigid pairs $(M,P)$ with
rigid reduced morphism objects $C_M\oplus Z_P$ and support $\tau$-tilting pairs
with maximal rigid objects after adjoining $K_A$ \cite{Li}.  This is the
categorical bridge used throughout the paper.

\subsection{Exact Hall algebra and normalization}

Assume now $\kk=\mathbb F_q$ and enlarge scalars to a field containing
$v=q^{1/2}$.  Let $\Hall(\MM)$ be the Ringel--Hall algebra with ordinary basis
$[X]$ indexed by isomorphism classes.  We use the extension-normalized basis
\[
 \delta_X=|\Aut X|[X]
\]
and the skew-normalized basis
\begin{equation}
 \eps_X=v^{-\langle X,X\rangle_{\MM}}\delta_X.
 \label{eq:skew-normalized-basis}
\end{equation}
The exact Euler form is
\[
 \langle X,Y\rangle_{\MM}
 =\dim\Hom_{\MM}(X,Y)-\dim\Ext^1_{\MM}(X,Y),
\]
and its antisymmetrization is
\[
 \Lambda_{\MM}(x,y)
 =\langle y,x\rangle_{\MM}-\langle x,y\rangle_{\MM}.
\]
If $\Ext^1_{\MM}(X,Y)=0$, then the only extension is split and
\begin{equation}
 \eps_X\diamond\eps_Y
 =v^{\Lambda_{\MM}([X],[Y])}\eps_{X\oplus Y}.
 \label{eq:split-Hall-product}
\end{equation}

The Grothendieck class of a morphism remembers only its projective terms:
\begin{equation}
 K_0(\MM)\simeq
 K_0^{\rm sp}(\proj A)\oplus K_0^{\rm sp}(\proj A),
 \qquad
 [P^-\to P^0]\longmapsto([P^-],[P^0]).
 \label{eq:term-lattice}
\end{equation}
This is the source of the differential-blindness of the canonical integration.
It is also exactly the lattice on which the correcting cocycle lives.

Let $L=K\oplus K$, where $K=K_0^{\rm sp}(\proj A)$, and let
\[
 D=(-I\ I):L\longrightarrow K,
 \qquad D(p,q)=q-p.
\]
Fix a Hall-compatible quantum coefficient system at the initial seed.  It
provides an integral skew form $S$ on the split-index lattice, or, with frozen
directions, an extended skew form whose mutable pullback is $S$.  The previous
paper proves that when $S$ is invariant under all presentation-to-cluster fan
transfers, the form
\begin{equation}
 \widehat\Lambda_S=D^TSD
 \label{eq:target-term-form}
\end{equation}
restricts to the quantum commutation form on every rigid chamber.  Define the
bicharacter defect
\begin{equation}
 \Omega_S=\widehat\Lambda_S-\Lambda_{\MM}
 \label{eq:Hall-cocycle}
\end{equation}
and twist homogeneous products by
\begin{equation}
 x\star_S y
 =v^{\Omega_S(\deg x,\deg y)}x\diamond y.
 \label{eq:twisted-Hall-product}
\end{equation}
Bilinearity of $\Omega_S$ makes $\star_S$ associative on the entire Hall
algebra.  If $X,Y$ belong to one rigid chamber, then
\begin{equation}
 \eps_X\star_S\eps_Y
 =v^{\Lambda_R(g_X,g_Y)}\eps_{X\oplus Y},
 \label{eq:compatible-Hall-quantum}
\end{equation}
which is exactly the based quantum-torus product in the corresponding seed.
For arbitrary rank, the canonical principal mutable parameter $S=0$ always
globalizes.  When $B_0$ is invertible and $B_0^T\Lambda_0=dI$, the parameter
$S=\Lambda_0$ gives the full-rank coefficient-free form.

\subsection{Central reduction and coefficient conventions}

Contractible terms should not produce new cluster variables.  Since
$D([K_P])=D([P],[P])=0$, the target form
\eqref{eq:target-term-form} kills every contractible direction.  Hence the
classes $\eps_{K_P}$ are central for $\star_S$.  We first localize at them and
then specialize them to the prescribed frozen boundary monomials.  Boundary
arcs are also inverted when we use the frozen-localized skein algebra.

\begin{definition}[Centrally reduced Hall algebra]
\label{def:Hred}
Let $\Hall_S^\star(\MM)$ be the cocycle-twisted Hall algebra.  The algebra
$\Hred$ is obtained by:
\begin{enumerate}[label=(\roman*),leftmargin=2.2em]
\item extending scalars to $\mathbb Q(\omega)$ and identifying the Hall
      parameter $v$ with the chosen power of $\omega$;
\item localizing at all $\eps_{K_{P_i}}$;
\item imposing the central coefficient relations which identify the
      contractible classes with the frozen boundary monomials of the chosen
      polygon coefficient system.
\end{enumerate}
\end{definition}

This operation is the analogue of the localization and central quotient in
acyclic morphism-Hall realizations.  It does not discard any reduced
non-rigid object.

\subsection{The polygon dictionary}

Assume from now on that $A$ is cluster-tilted of Dynkin type $A_n$.  Choose a
cluster category realization
\[
 A\simeq\End_{\CC}(T)
\]
and identify $\CC$ with the cluster category of the convex polygon
$P=P_{n+3}$.  The indecomposable objects of $\CC$ correspond to diagonals of
$P$; cluster-tilting objects correspond to triangulations; and
\[
 \Ext^1_{\CC}(X_\gamma,X_\delta)\ne0
 \quad\Longleftrightarrow\quad
 \gamma\text{ and }\delta\text{ cross}.
\]
Every indecomposable object of a Dynkin cluster category is rigid.

The functor from reduced projective morphisms to the two-term homotopy window,
followed by desuspension into $\CC$, identifies
\[
 C_M\quad\text{and}\quad Z_{P_i}
\]
with the polygon diagonals not in, respectively in, the initial triangulation.
The number of reduced indecomposable classes is
\[
 \frac{n(n+1)}2+n=\frac{n(n+3)}2,
\]
the number of diagonals of $P_{n+3}$.

For a reduced object
\[
 M=\bigoplus_{j=1}^r M_{\gamma_j}
\]
write $D(M)$ for the multidiagonal $\{\gamma_1,\ldots,\gamma_r\}$, with
multiplicities.  Define
\begin{equation}
 \crs(M)=
 \#\{(i,j):i<j,\ \gamma_i\text{ crosses }\gamma_j\}.
 \label{eq:crossing-number}
\end{equation}
Then $M$ is rigid if and only if $\crs(M)=0$.

\subsection{The whole reduced Hall algebra is generated by arcs}

The next elementary point is important: the source of the main theorem is the
whole centrally reduced Hall algebra.

\begin{lemma}[Indecomposable generation in finite representation type]
\label{lem:indecomposable-generation}
Let $\mathcal E$ be a finitary Krull--Schmidt exact category with finitely many
indecomposable isomorphism classes.  Over a coefficient field in which the
split Hall coefficients are invertible, $\Hall(\mathcal E)$ is generated by
the basis classes of indecomposable objects.
\end{lemma}

\begin{proof}
Let $M=M_1\oplus\cdots\oplus M_r$ be a Krull--Schmidt decomposition and
choose an order of its summands.  The Hall product
\[
 \delta_{M_1}\diamond\cdots\diamond\delta_{M_r}
\]
contains the split object $M$ with a nonzero coefficient.  Every other middle
term is a non-split iterated extension and degenerates to $M$.  Order
isomorphism classes by the degeneration order, refined by total dimension.
The displayed product is triangular with invertible diagonal coefficient.
Induction in that finite interval expresses $\delta_M$ in the algebra
generated by the indecomposable classes.  The same argument applies after a
bicharacter twist, localization and central specialization.
\end{proof}

Cluster-tilted algebras of Dynkin type are representation-finite, and
Bautista's normal form shows that $\Mor(\proj A)$ has only finitely many
indecomposable classes up to the contractibles.  Therefore
\cref{lem:indecomposable-generation} applies to $\Hred$.

\begin{corollary}[Arc generation]
\label{cor:arc-generation}
The algebra $\Hred$ is generated by the reduced indecomposable classes
$\eps_\gamma$ indexed by polygon diagonals together with the invertible frozen
boundary monomials.
\end{corollary}
